% GENERATED FILE. Edit canonical lessons, then run npm run generate:books.
% canonical-source-hash: fnv1a64:48f9d02f603cb114
% canonical-lessons: LA-C159-sentio, LA-C159-pendeo, LA-C159-taceo, LA-C159-servo, LA-C159-muto

\chapter{To feel, To hang, To be silent, To keep, To change}
\label{ch:la-to-feel-to-hang-to-be-silent-to-keep-to-change}
\hlchaptermodality{159}

\begin{chapteropening}
\textbf{By the end of this chapter:} \emph{I can say to feel, to hang, to be silent, to keep and to change.}

Complete the last lesson of chapter 159: I can say to feel, to hang, to be silent, to keep and to change.
\end{chapteropening}

\section[sentiō, sentīre]{sentiō, sentīre — to feel}

\label{lesson:LA-C159-sentio}



\begin{quote}
Before the new one: say the Latin for to dine, then the Latin for to taste.
\end{quote}

\subsection*{You'll want to know: sentiō, sentīre}
\textbf{sentiō, sentīre} — \textquotedblleft{}to feel\textquotedblright{}.

\begin{cousinweb}
\textbf{sentiō} is inherited, usually traced to Proto-Indo-European \emph{*sent-}, \textquotedblleft{}to head for, to go\textquotedblright{} (EtymDB 2.1). English \emph{send} is a cousin.
\end{cousinweb}

\begin{grammarlens}[title={What you've built}]
One more word for this chapter.
\end{grammarlens}

\subsection*{Your turn}
\begin{itemize}
\raggedright
  \item \emph{Say it:} \emph{sentiō}
  \item \emph{Say it:} \emph{sentiō}, once more
  \item \emph{Say it:} say \emph{sentiō}
  \item \emph{Recall:} say the Latin for to dine, then the Latin for to taste, then say \emph{sentiō} again
\end{itemize}

\begin{tcolorbox}[breakable,colback=teal!4,colframe=teal!35!black,title={Before you move on}]
What does sentiō, sentīre mean? (\textquotedblleft{}To feel\textquotedblright{}.) Say it once more.
\end{tcolorbox}
\section[pendeō, pendēre]{pendeō, pendēre — to hang}

\label{lesson:LA-C159-pendeo}



\begin{quote}
Before the new one: say the Latin for to taste, then the Latin for to feel.
\end{quote}

\subsection*{You'll want to know: pendeō, pendēre}
\textbf{pendeō, pendēre} — \textquotedblleft{}to hang\textquotedblright{}.

\begin{cousinweb}
\textbf{pendeō} is the intransitive partner of \emph{pendō}, \textquotedblleft{}to weigh, to hang up\textquotedblright{} (Lewis \& Short): to be hanging. The pair is usually traced to Proto-Indo-European \emph{*(s)pend-}, \textquotedblleft{}to stretch, to pull\textquotedblright{} (EtymDB 2.1).
\end{cousinweb}

\begin{grammarlens}[title={What you've built}]
One more word for this chapter.
\end{grammarlens}

\subsection*{Your turn}
\begin{itemize}
\raggedright
  \item \emph{Say it:} \emph{pendeō}
  \item \emph{Say it:} \emph{pendeō}, once more
  \item \emph{Say it:} say \emph{pendeō}
  \item \emph{Recall:} say the Latin for to taste, then the Latin for to feel, then say \emph{pendeō} again
\end{itemize}

\begin{tcolorbox}[breakable,colback=teal!4,colframe=teal!35!black,title={Before you move on}]
What does pendeō, pendēre mean? (\textquotedblleft{}To hang\textquotedblright{}.) Say it once more.
\end{tcolorbox}
\section[taceō, tacēre]{taceō, tacēre — to be silent}

\label{lesson:LA-C159-taceo}



\begin{quote}
Before the new one: say the Latin for to feel, then the Latin for to hang.
\end{quote}

\subsection*{You'll want to know: taceō, tacēre}
\textbf{taceō, tacēre} — \textquotedblleft{}to be silent\textquotedblright{}.

\begin{cousinweb}
\textbf{taceō} is inherited from Proto-Italic \emph{*takēō}, \textquotedblleft{}to be silent\textquotedblright{} (EtymDB 2.1), but Lewis \& Short call its deeper origin doubtful. Gothic \emph{þahan}, \textquotedblleft{}to be silent\textquotedblright{}, is often compared.
\end{cousinweb}

\begin{grammarlens}[title={What you've built}]
One more word for this chapter.
\end{grammarlens}

\subsection*{Your turn}
\begin{itemize}
\raggedright
  \item \emph{Say it:} \emph{taceō}
  \item \emph{Say it:} \emph{taceō}, once more
  \item \emph{Say it:} say \emph{taceō}
  \item \emph{Recall:} say the Latin for to feel, then the Latin for to hang, then say \emph{taceō} again
\end{itemize}

\begin{tcolorbox}[breakable,colback=teal!4,colframe=teal!35!black,title={Before you move on}]
What does taceō, tacēre mean? (\textquotedblleft{}To be silent\textquotedblright{}.) Say it once more.
\end{tcolorbox}
\section[servō, servāre]{servō, servāre — to keep, to save}

\label{lesson:LA-C159-servo}



\begin{quote}
Before the new one: say the Latin for to hang, then the Latin for to be silent.
\end{quote}

\subsection*{You'll want to know: servō, servāre}
\textbf{servō, servāre} — \textquotedblleft{}to keep, to save\textquotedblright{}.

\begin{cousinweb}
\textbf{servō} is inherited, usually traced to Proto-Indo-European \emph{*ser-}, \textquotedblleft{}to watch over, to protect\textquotedblright{} (EtymDB 2.1). It is probably not the same word as \emph{servus}, \textquotedblleft{}a slave\textquotedblright{}.
\end{cousinweb}

\begin{grammarlens}[title={What you've built}]
One more word for this chapter.
\end{grammarlens}

\subsection*{Your turn}
\begin{itemize}
\raggedright
  \item \emph{Say it:} \emph{servō}
  \item \emph{Say it:} \emph{servō}, once more
  \item \emph{Say it:} say \emph{servō}
  \item \emph{Recall:} say the Latin for to hang, then the Latin for to be silent, then say \emph{servō} again
\end{itemize}

\begin{tcolorbox}[breakable,colback=teal!4,colframe=teal!35!black,title={Before you move on}]
What does servō, servāre mean? (\textquotedblleft{}To keep, to save\textquotedblright{}.) Say it once more.
\end{tcolorbox}
\section[mūtō, mūtāre]{mūtō, mūtāre — to change}

\label{lesson:LA-C159-muto}



\begin{quote}
Before the new one: say the Latin for to be silent, then the Latin for to keep.
\end{quote}

\subsection*{You'll want to know: mūtō, mūtāre}
\textbf{mūtō, mūtāre} — \textquotedblleft{}to change\textquotedblright{}.

\begin{cousinweb}
\textbf{mūtō} is usually traced to Proto-Indo-European \emph{*mey-}, \textquotedblleft{}to change, to exchange\textquotedblright{}, the root of \emph{mūtuus}, \textquotedblleft{}borrowed, mutual\textquotedblright{}; Lewis \& Short's link with \emph{moveō}, \textquotedblleft{}to move\textquotedblright{}, is less likely.
\end{cousinweb}

\begin{grammarlens}[title={What you've built}]
One more word for this chapter.
\end{grammarlens}

\subsection*{Your turn}
\begin{itemize}
\raggedright
  \item \emph{Say it:} \emph{mūtō}
  \item \emph{Say it:} \emph{mūtō}, once more
  \item \emph{Say it:} say \emph{mūtō}
  \item \emph{Recall:} say the Latin for to be silent, then the Latin for to keep, then say \emph{mūtō} again
\end{itemize}

\begin{tcolorbox}[breakable,colback=teal!4,colframe=teal!35!black,title={Before you move on}]
What does mūtō, mūtāre mean? (\textquotedblleft{}To change\textquotedblright{}.) Say it once more.
\end{tcolorbox}
